\section{Introduction}
\label{sec:intro}

In his 2009 essay \emph{Birds and Frogs}, Freeman Dyson proposed a way of
looking at the Riemann Hypothesis that is geometric rather than analytic
\cite{dyson2009}. If RH is true, the nontrivial zeros of the zeta function form
a \emph{one-dimensional quasicrystal}: a point set that is not periodic, yet
whose diffraction spectrum is pure point---discrete spectral lines, like a
crystal's---supported on the logarithms of the prime powers. Classify the
one-dimensional quasicrystals, Dyson suggested, and one might find the zeros on
the list. The mathematical substance beneath the metaphor is not conjectural at
all: it is the Guinand--Weil explicit formula \cite{guinand1948,weil1952}, the
identity under which weighted sums over the zeros equal weighted sums over the
primes---the zeros \emph{diffracting} off the primes. Everything in the
quasicrystal picture descends from that identity, which in turn descends from
Riemann's 1859 memoir.

This paper reports a body of machine-checked mathematics---formalized in
Lean~4 and re-proved from scratch by Mathlib's kernel---that carries the
\emph{finite} part of that program, and is scrupulous about the boundary
between what is proved and what is not. Three results form the spine
(\cref{sec:rh}):
\begin{itemize}
  \item \textbf{A finite-height explicit formula, in the kernel.} For any box
    with right edge in $\operatorname{Re} s > 1$ and any holomorphic test
    weight, weighted sums over zeta's \emph{actual} meromorphic divisor equal
    von Mangoldt prime-power sums plus three explicit boundary integrals---a
    rectangular Guinand--Weil identity, the diffraction identity at finite
    height. The prime-side edge needs no numeric certificate at all: its
    non-vanishing is a theorem.
  \item \textbf{All zeros to height 4000, on the line.} A complete finite
    Turing verification: the $3474$ nontrivial zeros with
    $0 < \operatorname{Im}\rho \le 4000$ lie exactly on
    $\operatorname{Re} s = \tfrac12$, assembled from eighty-nine band
    certificates composed through an unbounded-depth height chain, their
    winding counts derived from boundary contours, with
    a uniform, documented numeric trust boundary and no bare classical
    hypothesis.
  \item \textbf{An effective de la Vallée-Poussin region.} A self-contained,
    machine-checked $\zeta(s) \ne 0$ on
    $\operatorname{Re} s > 1 - c/\log|\operatorname{Im} s|$, with $c$ an
    explicit closed-form constant---a refinement over the classically
    non-effective statement---together with elementary Hadamard-free regions
    that appear to be the first of their kind verified in Lean.
\end{itemize}
Around that spine the same section reports, at the same standard of honesty,
finite kernel-checked prefixes of two classical RH criteria (Jensen--P\'olya
hyperbolicity and Li positivity) and---asserted by Riemann in 1859, proved by
von Mangoldt, and now a theorem of the kernel---the Riemann--von Mangoldt
counting formula itself, $N(T) = \theta(T)/\pi + 1 + S(T)$, exact at finite
height.
Let us say immediately what none of this is: \emph{we do not claim the Riemann
Hypothesis}. The explicit formula is an instrument for studying RH, not a
conclusion about it; the Turing verification is finite; the zero-free region is
a region. Every numeric input enters as an explicit, documented hypothesis, and
\texttt{conjecture1\_proved = False} is a real, tested invariant of the
repository, not a slogan.

\paragraph{The instrument.} Results at this scale---hundreds of kernel-checked
theorems, many of them generated rather than hand-written---raise a trust
question. Search procedures, heuristics, and increasingly large language models
are very good at producing arguments that \emph{look} right, and the more
fluent the generator, the more convincing its mistakes. Our answer is old and
clean: \emph{check the witness, don't trust the author}. The instrument behind
every result above is \emph{Telperion}, a certificate compiler for Lean~4: you
describe a family of concrete claims, Telperion certifies each instance in
exact rational arithmetic and emits Lean, and Mathlib's kernel re-proves every
line from scratch (\cref{sec:system}). The generator is untrusted by design---a
wrong certificate is a compile error, never a false theorem. Because the kernel
cannot catch a theorem that is true but \emph{vacuous}, Telperion adds a
self-verification layer---per-emitter sensitivity declarations,
forge-the-witness negative controls, and an independent second checker---and is
honest about the one thing no machinery can settle: whether a formal statement
means the informal claim a human intended (\cref{sec:vacuity}).

\paragraph{The instrument elsewhere.} The same discipline carries to problems
far from the zeta function, and we include two of them as evidence that the
method, not the subject, is what generalizes: a machine-checked analytic core
for the (still open) Brualdi--Goldwasser tree-maximizer conjecture, where we
say exactly which bridge is missing (\cref{sec:bg}); and a kernel-checked,
symbolic-in-$n$ sum-of-squares lower bound for the Grigoriev knapsack system in
proof complexity, where the mathematics is due to others and our contribution
is the certified pipeline (\cref{sec:pc}).

\paragraph{A note on posture.} A reader will notice that we spend as much care
on what is \emph{not} proved as on what is. That is deliberate. The honesty
flags in this project are load-bearing, and a certificate engine whose whole
value is that you need not trust it would be a strange place to start
overclaiming. \Cref{sec:evaluation} takes up the limitations directly, and the
reproducibility appendix gives a reviewer the commands to re-check any claim in
this paper themselves.
